%--------------------------------------------------------------------------------
%
%
%--------------------------------------------------------------------------------
\pagebreak
\section*{WN — Create a Window}
\addcontentsline{toc}{section}{WN — Create a Window}

%--------------------------------------------------------------------------------
\begin{iPageDescEntry}{Syntax}
\texttt{wn <type> [ , <arg1> [ , <arg2> ]] } \\
\end{iPageDescEntry}

%--------------------------------------------------------------------------------
\begin{iPageDescEntry}{Description}

The \texttt{WN} command creates a new window. The mandatory \texttt{<type>}
parameter specifies the type of window to create. Optional arguments provide
additional data required for the window creation.

The newly created window becomes the current window and is initialized with
default attributes, such as the number of lines, layout, and numeric radix.

\begin{flushleft}
\begin{tabularx}{\textwidth}{|l|X|X|p{5cm}|}
\hline
\textbf{Type} & \textbf{Arg 1} & \textbf{Arg 2} & \textbf{Action} \\ \hline
\texttt{MEM}         & \texttt{memAdr} & \texttt{-} & Create a memory window. \\ \hline
\texttt{PROC}        & \texttt{modNum} & \texttt{-} & Create all PROC-related windows (CPU, ITLB, DTLB, ICACHE, DCACHE) for the module. \\ \hline
\texttt{CPU}         & \texttt{modNum} & \texttt{-} & Create a CPU window. \\ \hline
\texttt{ITLB}        & \texttt{modNum} & \texttt{-} & Create an instruction TLB window. \\ \hline
\texttt{DTLB}        & \texttt{modNum} & \texttt{-} & Create a data TLB window. \\ \hline
\texttt{ICACHE}      & \texttt{modNum} & \texttt{-} & Create an instruction cache window. \\ \hline
\texttt{DCACHE}      & \texttt{modNum} & \texttt{-} & Create a data cache window. \\ \hline
\texttt{TEXT}        & \texttt{filePath} & \texttt{-} & Create a text window displaying a file. \\ \hline
\texttt{CODE}        & \texttt{memAdr} & \texttt{-} & Create a memory window showing assembly code. \\ \hline
\end{tabularx}
\end{flushleft}

The window type \texttt{PROC} is a collection that includes CPU, ITLB, DTLB,
ICACHE, and DCACHE windows for the specified module.

\end{iPageDescEntry}

%--------------------------------------------------------------------------------
\begin{iPageDescOpEntry}{Example}
\begin{lstlisting}[style=iPageOpStyle]
wn mem, 0x1000      # Create a memory window at address 0x1000.
wn cpu, 3           # Create a CPU window for module 3.
wn dtlb, 4          # Create a data TLB window for module 4.
\end{lstlisting}
\end{iPageDescOpEntry}

%--------------------------------------------------------------------------------
\begin{iPageDescEntry}{Notes}
All newly created windows start with default attributes, which can be modified
later using commands like \texttt{WL} (set window lines) or \texttt{WR} (set window radix), if applicable for the window.
\end{iPageDescEntry}
